% Full English proofs, independently reorganized from the 2026-10-14 record.
% This is an input file, not a standalone document.
\section{Probability tools and reporting conventions}
\label{app:tools}

For the reporting constructions, index arms in arrival order
$i=1,2,\ldots$, with a fixed order within each batch.
Every arm has an independent potential reward stream $(Y_{i,n})_{n\ge1}$
in $[0,1]$. The policy observes its $n$th entry on its $n$th pull. Arrival
times are deterministic, and reports are issued after the current observation.
Write $N_i(t)$ for lifetime observations through round $t$ and
$\hat r_i(n)=n^{-1}\sum_{s\le n}Y_{i,s}$. The observed reward in the
original physical-time notation is
$X_{t,I_t}=Y_{I_t,N_{I_t}(t-1)+1}$. The conclusions below also
hold conditionally on an independent exogenous arrival schedule.

For independent $[0,1]$ observations, Hoeffding's inequality gives
\[
 P\{|\hat r_i(n)-r_i|>x\}\le2e^{-2nx^2}.
\]
For completeness, the centered log moment generating function has value and
derivative zero at the origin; its second derivative is a variance under an
exponentially tilted distribution and is at most $1/4$. Integrating twice
gives $Ee^{\lambda(X-EX)}\le e^{\lambda^2/8}$. Multiplication over independent
observations and exponential Markov inequality prove the display.

With
\[
 b_i(n)=\sqrt{\frac{\log(8i^2n^2/\delta)}{2n}},
\]
the union of all potential-prefix failures has probability at most
$\delta\zeta(2)^2/4<\delta$. Intersecting the prefix intervals with $[0,1]$
therefore gives simultaneous coverage. On that event each endpoint differs
from its mean by at most $2b_i(n)$. No conditioning on a previously selected
empirical mean is used. Empty intersections on failure paths are handled by
exclusion-first reporting; the reward policy remains defined on all paths.

The checkpoint alternative fixes, before observing arm $i$,
\[
 q_i(\rho_i)=\left\lceil\frac2{\rho_i^2}
            \log\frac{4i(i+1)}\delta\right\rceil.
\]
Before $q_i$ observations its interval is $[0,1]$; afterwards it is permanently
\[
 [\hat r_i(q_i)-\rho_i/2,
  \hat r_i(q_i)+\rho_i/2]\cap[0,1].
\]
The sum of coverage failures is at most
$\sum_i\delta/\{2i(i+1)\}=\delta/2$. On coverage, each endpoint differs from
the true mean by at most $\rho_i$. These intervals are nonempty on every path
and only contract once. If $\underline r^*(t)=\max_i\ell_i(t)$ and
$\overline r^*(t)=\max_i u_i(t)$, use
\[
 \mathcal P_t=\{i:\ell_i(t)\ge\overline r^*(t)-\varepsilon\},\quad
 \mathcal N_t=\{i:u_i(t)<\underline r^*(t)-\varepsilon\},\quad
 \mathcal V_t=\mathcal{A}_t\setminus(\mathcal P_t\cup \mathcal N_t).
\]
The sets $\mathcal P_t,\mathcal N_t$ are disjoint because
$\underline r^*(t)\le\overline r^*(t)$ and $\ell_i\le u_i$.
Coverage implies
$\mathcal P_t\subseteq \mathcal G_t^\varepsilon\subseteq \mathcal{A}_t\setminus \mathcal N_t$.
Moreover, $\mathcal N_t$ is increasing as observations and arrivals accumulate:
$\underline r^*(t)$ cannot decrease and an existing $u_i$ cannot increase.

For two Bernoulli mean vectors $\boldsymbol r=(r_i)_i$ and
$\boldsymbol r'=(r_i')_i$ and a bounded stopping time $\tau$, action
kernels, policy randomization, and stopping kernels cancel in the history
likelihood ratio. Padding after stopping by null symbols and taking iterated
expectations yields
\begin{equation}
 D(P_{\boldsymbol r}^{H_\tau}\Vert P_{\boldsymbol r'}^{H_\tau})
 =\sum_iE_{\boldsymbol r} N_i(\tau)\operatorname{kl}(r_i,r_i').
 \label{eq:app_kl_identity}
\end{equation}
The same proof applies to arbitrary observation laws with finite two-way KL.
For an event $F$ measurable at stopping, splitting the likelihood integral
over $F,F^c$ and applying Jensen gives
\[
 D(P\Vert Q)\ge\operatorname{kl}(P(F),Q(F)).
\]
We repeatedly use
\[
 2(p-q)^2\le\operatorname{kl}(p,q)
 \le\frac{(p-q)^2}{q(1-q)}.
\]
The upper bound follows from $\log x\le x-1$; for the lower bound the second
derivative in $p$ is $1/\{p(1-p)\}\ge4$ and both value and derivative vanish
at $p=q$.

\section{The broadcast obstruction and its queue consequence}
\label{app:broadcast}

\paragraph{Precise claim.}
Let $0<\varepsilon\le1/16$, $0<\delta\le1/16$,
$0<\eta<\varepsilon/2$. The reference arm has mean
$1/4+\varepsilon/2$. All candidates have mean $p=1/4$ in environment zero;
environment $j$ raises only candidate $j$ to $p+2\varepsilon$. Let $M_t$ be
the candidate count, $a_t$ the new-candidate count, and
\[
 q_- =\left\lfloor\frac{\operatorname{kl}(1/4,\delta)}
                         {\operatorname{kl}(1/4,1/4+2\varepsilon)}\right\rfloor,
 \quad Q_t=(Q_{t-1}+q_-a_t-1)_+,\quad Q_0=0,
\]
\[
 \mathcal B_T=\sum_{t\le T}(M_t-1)_+1\{Q_t>0\}.
\]
Every policy reliable on these environments obeys
\[
 E_0\mathcal U_T(\eta)\ge(3/4-\delta)\mathcal B_T,
 \qquad
 P_0\{\mathcal U_T(\eta)\ge\mathcal B_T/4\}\ge1/4
\]
when $\mathcal B_T>0$.

\begin{proof}
The denominator defining $q_-$ is at most
$\operatorname{kl}(1/4,3/8)\le1/15$, whereas its numerator is at least
$\frac14\log4+\frac34\log(4/5)\ge7/48>2/15$. Thus $q_-\ge2$.
The recursion gives, by induction,
\[
 Q_t=\max_{0\le u\le t}\left\{
       q_-\sum_{r=u+1}^ta_r-(t-u)\right\}.
\]
If $Q_t>0$, there is a suffix in which total candidate demand exceeds all
physical observation slots. Its arms had no observations before that suffix.
Consequently at least one deterministic candidate label $j$ satisfies
$E_0N_j(t)<q_-$. The choice is made from expected counts, not from a realized
history. For any other candidate $i$, membership of $i$ in $\mathcal P_t$ is erroneous
under environment $j$. Its probability there is at most $\delta$.
The history KL is strictly below $\operatorname{kl}(1/4,\delta)$, so data
processing and monotonicity of binary KL force $P_0(i\in \mathcal P_t)<1/4$.
Since all candidates are genuinely good under environment zero,
$P_0(i\in \mathcal N_t)\le\delta$. It follows that
$P_0(i\in \mathcal V_t)\ge3/4-\delta$. These candidates all have boundary distance
$\varepsilon/2>\eta$. Summing at deterministic busy rounds proves the
expectation bound, entirely under the one environment zero.

Let $Y$ count candidate unresolved area only at busy rounds with $M_t\ge2$.
Then $Y\le\mathcal U_T$, $0\le Y\le2\mathcal B_T$, and
$E_0Y\ge11\mathcal B_T/16$. If
$p=P_0(Y\ge\mathcal B_T/4)$, then
$11\mathcal B_T/16\le\mathcal B_T/4+(7\mathcal B_T/4)p$, so $p\ge1/4$.
\end{proof}

\paragraph{Queue area.}
For an integer $q\ge2$, suppose $m$ candidates arrive by time $s$ and put
$D=s+qm-1$. Then
\[
 \sum_{t\le D}(M_t-1)_+1\{Q_t(q)>0\}
 \ge\frac{q-1}{2}m(m-1).
\]
To prove this, remove all subsequent arrivals, which can only decrease both
the queue and its weights. The $qm$ service units finish by $D$. At the
$n$th actual service, at least $\lceil n/q\rceil$ candidates have arrived,
so service rounds weighted by $M_t-1$ sum to at least
$\sum_{n=1}^{qm}(\lceil n/q\rceil-1)=qm(m-1)/2$.
The post-service queue is zero at the last service round of each busy period.
Their cumulative arrival counts are distinct integers in $\{1,\ldots,m\}$,
so removing their weights costs at most $m(m-1)/2$. This proves the claim.

In particular, if along a fixed infinite schedule a policy reliable in every
environment has a deterministic area envelope $a_T=o(T)$, then
$K_t=o(\sqrt t)$. Otherwise along $s_n$ there are at least $c\sqrt{s_n}$
candidates. Choose the first $m_n=\lfloor c\sqrt{s_n}\rfloor$; eventually
$q_-m_n\le s_n$. The queue bound and the preceding probability statement at
$2s_n$ give $P_0(\mathcal U_{2s_n}\ge c's_n)\ge1/4$, contradicting
$\delta\le1/16$ and the proposed envelope.

\section{The partial simultaneous-recovery lemma}
\label{app:partial}

\paragraph{Precise claim.}
Let $m\ge1$ independent uniform signs specify independent observation laws
$P_i^+,P_i^-$. All nuisance parameters and observations are independent of
these signs. Sampling may be adaptive and interleaved, and $\tau$ is bounded.
At every sign configuration, the probability of any erroneous determinate
report is at most $\delta\le1/16$. Let $J_\tau$ be the number of determinate
reports and
$\kappa=\max_i\max\{D(P_i^+\Vert P_i^-),D(P_i^-\Vert P_i^+)\}\in(0,\infty)$.
If $\overline P(J_\tau\ge m/2)\ge1/2$, then
\[
 \overline E\sum_iN_i(\tau)\ge\frac{m}{1024\kappa}\log\frac m\delta.
\]

\begin{proof}
Let $L_i$ be the observed positive-versus-negative log likelihood ratio for
coordinate $i$. After conditioning on the full stopped history, the action
and stopping kernels cancel. The likelihood factors over coordinates and so
does the posterior, even though the sample counts need not be independent.
The smallest coordinatewise posterior error is
$e_i=(1+e^{|L_i|})^{-1}$. If $S$ is the reported subset, its conditional
probability of complete correctness is at most $\prod_{i\in S}(1-e_i)$;
choosing a non-MAP label cannot improve it. Thus
\[
 \overline E\left[1-\prod_{i\in S}(1-e_i)\right]\le\delta.
\]
With probability at least $1-2\delta$ this product is at least $1/2$ and hence
$\sum_{i\in S}e_i\le\log2$. Intersecting with $J_\tau\ge m/2$ costs at most
$2\delta$ and leaves probability at least $3/8$. For $m\ge4096$, Jensen gives
on this intersection
\[
 \sum_i|L_i|\ge\sum_{i\in S}\{\log(1/e_i)-\log2\}
 \ge\frac m2\log\frac{m}{4\log2}.
\]
Successive conditional expectations of signed likelihood increments give
\[
 \overline E\sum_i\sigma_iL_i\le\kappa\overline E\sum_iN_i(\tau).
\]
Finite two-way KL makes the stopped increments integrable. Given the history,
$E(\sigma_iL_i\mid H_\tau)=|L_i|(1-2e_i)$, and
$2|L_i|e_i\le1$ for every real $L_i$. Therefore
\[
 \kappa\overline E\sum_iN_i\ge\overline E\sum_i|L_i|-m
 \ge m\left\{\frac3{16}\log\frac{m}{4\log2}-1\right\}
 \ge\frac m{256}\log m,
\]
where the last inequality holds for $m\ge4096$.

For the confidence term, fix a sign vector and flip coordinate $i$ only.
The event reporting the original sign on coordinate $i$ has original
probability at least its determinate-report probability minus $\delta$, and
flipped probability at most $\delta$. Average the data-processing inequalities
over vectors and coordinates and use joint convexity of binary KL.
Since $\overline EJ_\tau\ge m/4$, this yields
\[
 \kappa\overline E\sum_iN_i\ge
 m\operatorname{kl}(1/4-\delta,\delta)
 \ge\frac m{256}\log(1/\delta).
\]
To check the last scalar inequality, write $L=\log(1/\delta)$. For $L\le8$,
Pinsker gives a lower bound $2(1/4-2\delta)^2\ge1/32\ge L/256$.
For $L\ge8$, binary entropy at most $\log2$ gives
$\operatorname{kl}(1/4-\delta,\delta)\ge3L/16-\log2\ge L/16$.
When $m\ge4096$, the maximum of the two preceding bounds is at least half
their sum. When $m<4096$, $\log m\le3\log(1/\delta)$, and the confidence
bound alone suffices. Both cases imply the stated constant $1/1024$.
\end{proof}

\paragraph{Single-window truncation.}
If such a batch has $H$ rounds with qualifying boundary distance larger than
$\eta$, set
\[
 h_* =\min\left\{H,
 \left\lceil\frac{m\log(m/\delta)}{1024\kappa}\right\rceil-1\right\}.
\]
At every $r\le h_*$ there are at most $r$ batch observations. The recovery
lemma's contrapositive makes the probability of half-completion less than
$1/2$, so the expected unresolved count is at least $m/4$. Thus batch area
has expectation at least $mh_*/4$. Its first-$h_*$-round area $Y$ lies in
$[0,mh_*]$; consequently
$\overline P(Y\ge mh_*/8)\ge1/7$. One fixed sign configuration attains
at least this probability. If $h_*=0$, only the trivial zero lower bound is
asserted.

\section{Area-constrained multiple updates}
\label{app:area_lower}

In window $e$, suppose $M_e$ candidates have gaps
$\varepsilon-\sigma_{e,k}d_e$, and a new reference makes their signs
irrelevant after that window. Put $w_e=M_eH_e$, $V=\sum_ew_e$, and let
$p_e$ be the product-prior probability that at least half of batch $e$ becomes
determinate simultaneously before its window ends. Define
\[
 \Phi_\delta(u)=
 \begin{cases}
  \operatorname{kl}(1-\delta-u,\delta),&u<1-2\delta,\\
  0,&u\ge1-2\delta,
 \end{cases}
\]
\[
 g_{m,\delta}(p)=\max\left\{0,
 \frac{(p-2\delta)_+}{2}
       \left[\log\frac m{4\log2}\right]_+-1,
 \Phi_\delta(1-p/2)\right\}.
\]
Each displayed component is convex and nondecreasing in $p$.

\begin{proof}[Proof of the general area-to-regret lower bound]
If window $e$ never half-completes, at every round more than half its
candidates remain unresolved. Its area is at least $w_e/2$. On area success,
$\sum_ew_e1_{S_e^c}\le2a$; outside that event this sum is at most $V$.
Consequently
\[
 \sum_ew_ep_e\ge(1-\delta)V-2a.
\]
Stop batch $e$ at first half-completion or at its endpoint. The likelihood
argument above now intersects the product-posterior event with completion
probability $p_e$, giving intersection probability at least
$(p_e-2\delta)_+$ and expected determinate count at least $M_ep_e/2$.
The absolute-likelihood and coordinate-flip derivations therefore give
\[
 \kappa_e\overline EN_e\ge M_eg_{M_e,\delta}(p_e).
\]
For the absolute-likelihood term, when the bracketed logarithm is positive
its Jensen expression is increasing for subset sizes at least $M_e/2$;
when it is nonpositive the stated term is at most zero. Each observation
counted before stopping has regret at least $\varepsilon-d_e$.
The signs of all batches belong to one product prior, and samples of
different batches are disjoint. Hence
\[
 \sup_\sigma E_\sigma R_T\ge
 \inf_{\substack{0\le p_e\le1\\
       \sum_ew_ep_e\ge(1-\delta)V-2a}}
 \sum_e\frac{(\varepsilon-d_e)M_e}{\kappa_e}
                  g_{M_e,\delta}(p_e).
\]
This is an information lower relaxation, not a claim that every feasible
allocation is achieved by a causal policy.
\end{proof}

\paragraph{Homogeneous consequence.}
Take $M_e=M$, $H_e=H$, $d_e=d$, and $a=JMH/8$. On success at least $3J/4$
windows half-complete; thus $\sum_ep_e\ge45J/64$. At least $3J/8$ windows
have $p_e\ge1/2$: otherwise the sum is at most
$3J/8+(5J/8)/2=11J/16<45J/64$.
The Bernoulli means lie in $[1/4,3/4]$, giving $\kappa_e\le32d^2$.
The partial-recovery lemma and $\varepsilon-d\ge\varepsilon/2$ yield
\[
 \sup_\sigma E_\sigma R_T\ge
 \frac{3J}{8}\frac{\varepsilon M}{2\cdot1024\cdot32d^2}
             \log\frac M\delta
 \ge2^{-20}\frac{\varepsilon JM}{d^2}\log\frac M\delta.
\]

\section{History-retaining checkpoint UCB and evidence reuse}
\label{app:reuse}

CP-UCB-AA fixes each arm's checkpoint before seeing its data. At even
physical rounds it samples the earliest-arriving arm below its lifetime
checkpoint, if one exists. At all other rounds it maximizes, over every
available arm including reported exclusions,
\[
 J_i(t)=\hat r_i(N_i(t-1))+
             \sqrt{3\log(2t)/N_i(t-1)},
\]
with index $+\infty$ at zero count and deterministic ties. Every observation
serves both the lifetime checkpoint and the empirical reward estimate.

\paragraph{A reward bound that does not condition on report coverage.}
At physical time $t$, order arms by their first observed rewards. At most
$t$ streams have been exposed, each at at most $t$ potential prefixes. A
previously unobserved stream is independent of the history selecting it.
Applying Hoeffding to each newly exposed stream and union bounding gives
\[
 P\left\{\exists\text{ observed arm and prefix }n\le t:
 |\hat r_i(n)-r_i|>
        \sqrt{3\log(2t)/n}\right\}
 \le2t^2(2t)^{-6}=\frac1{32t^4}.
\]
The exposure order is used only for this union bound, not to assert
independence among observed counts. An ordinary UCB pull with finite index,
on this event's complement, has
$\Delta_{I_t}(t)\le2\sqrt{3\log(2t)/N_{I_t}(t-1)}$.
Each arm--pre-pull-count pair is charged at most once; interleaved service
observations skip counts. Therefore the ordinary, nonfirst-pull regret is
at most
\[
 4\sqrt{3\log(2T)}\sum_i\sqrt{N_i(T)}
 \le4\sqrt{3K_TT\log(2T)}
\]
on good rounds. First pulls cost at most $K_T$, and the sum of expected
bad-round losses is less than one. Adding the actual regret of service
rounds proves
\[
 ER_T\le\min\left\{T,K_T+4\sqrt{3K_TT\log(2T)}+1+
       E\sum_{t\in\mathcal I_T}\Delta_{I_t}(t)\right\}.
\]
There is no $\delta T$ term: the reward policy never permanently discards an
arm because of a reporting error.

\paragraph{Arbitrary-arrival completion bound.}
At any fixed reporting precision $\rho=\eta/4$, if all current checkpoints
are complete, all endpoint errors are at most $\eta/4$ and every arm whose
boundary distance exceeds $\eta$ has a determinate report. A virtual FIFO
queue adding $q_i$ work at each arrival and serving one unit at every even
round dominates unfinished checkpoints; additional UCB observations only
reduce actual work. Its busy periods, including alignment to the first even
round, have total length at most $2\sum_iq_i$. Thus on simultaneous coverage,
\[
 \mathcal U_T(\eta)\le
 \sum_{t\le T}K_t1\{\overline Q_t>0\}
 \le\min\{K_TT,2K_T\sum_{i\le K_T}q_i\}.
\]
At fixed $\eta,\delta$, the checkpoints are $O(\log(ei))$, so
$K_T^2\log(eK_T)=o(T)$ suffices for both sublinear area and expected regret.
Combined with the broadcast consequence, a path with
$K_t\asymp t^\alpha$, $0<\alpha<1$, admits both if and only if $\alpha<1/2$.
This is an exponent statement and does not remove the critical logarithm.

\paragraph{Separated windows and reuse.}
Let batch $e$ have reference mean $b_e$, $M_e$ candidate means
$b_e-\varepsilon\pm d_e$, and common predetermined precision $\rho_e$.
The final batch contains only its reference; set $M_J=M_{-1}=0$ and
$S_e=\sum_{i\in F_e}q_i(\rho_e)$. Assume
\[
 H_e\ge2S_e,\quad b_J-b_0\le\varepsilon/4,\quad
 d_e+\rho_e+\rho_{e+1}<b_{e+1}-b_e,\quad
 2\max_e\rho_e\le\varepsilon/2.
\]
The algorithm is not told which arm is the reference or when the next window
starts. By induction every batch completes its checkpoints within its first
$2S_e$ rounds, before the next batch appears. On coverage, successive
reference levels satisfy $b_e+\rho_e<b_{e+1}-\rho_{e+1}$. Once batch $e$
completes, the maximum upper endpoint of all available arms is at most
$b_e+\rho_e$. Its candidates are therefore completely classified when
$2\rho_e<d_e$.

An old candidate in batch $e-1$ has upper endpoint at most
$b_{e-1}-\varepsilon+d_{e-1}+\rho_{e-1}$, whereas the new reference has
lower endpoint at least $b_e-\rho_e$. Their difference exceeds
$\varepsilon$, excluding the entire preceding candidate batch without any
further old-arm observations. Older exclusions persist. Every reference has
gap at most $\varepsilon/4$ and total endpoint error at most
$\varepsilon/2$, so all references become positively certified.
If $F=\{e<J:2\rho_e<d_e\}$, the area is consequently at most
\[
 \sum_{e\notin F}M_eH_e+2\sum_{e\in F}M_eS_e+
                2\sum_{e=0}^J(e+1+M_{e-1})S_e.
\]
The terms cover current nonfine candidates, fine candidates, and temporarily
unresolved references plus the preceding candidate batch. All earlier
candidate batches have already been excluded.

For expected regret, every observation with pre-pull count below $q_i$
occurs in its arm's birth window. Such reference observations have zero
regret, and candidate observations cost at most $\varepsilon+d_e$ each.
All observations with pre-pull count at least $q_i$ are ordinary UCB pulls.
Let $\gamma_i$ be the smallest positive gap arm $i$ ever has and
$D_i^T$ its largest gap through $T$. At a good UCB round, any positive-gap
pull has pre-pull count at most $12\log(2T)/\gamma_i^2$. Since counts do not
repeat, this gives the unconditional bound
\[
 ER_T\le\sum_{e<J}\sum_{k=1}^{M_e}(\varepsilon+d_e)q_{e,k}
 +\sum_{i:\gamma_i\text{ exists}}D_i^T
 \left[\left\lceil\frac{12\log(2T)}{\gamma_i^2}\right\rceil-q_i+1\right]_+
 +1.
\]
The completion-before-next-batch assumption is essential to the first
charging term. It is not an assertion about arbitrary overlapping arrivals.

\section{The rapid and deferred endpoints on one family}
\label{app:tradeoff}

Here $J\ge1$, $M\ge\max\{2,J\}$,
$0<\varepsilon\le1/8$, $0<\delta\le1/16$, $Jv\le\varepsilon/4$,
$0<d\le v/32$, and $\eta=d/2$. The reference means are
$b_e=1/2+ev$ and the $M$ candidates born in window $e<J$ have means
$b_e-\varepsilon+\sigma_{e,k}d$. Each window has length $H$, and the final
window contains only a new reference. Set
\[
 K=J(M+1)+1,\quad L_K=\log\frac{4K(K+1)}\delta,
 \quad q_f=\lceil32L_K/d^2\rceil,\quad q_c=\lceil128L_K/v^2\rceil,
\]
\[
 H=128(M+J+1)q_f,\qquad T=(J+1)H,\qquad V=JMH.
\]
These are analysis quantities: the policy uses the selected precision
$\rho_f=d/4$ or $\rho_c=v/8$, not $K,H,T$ or the best-arm identities.
Both choices satisfy the separated-window conditions.

The preceding area bound is at most $B_{J,M}q_f$ for the fine policy and
$V+B_{J,M}q_c$ for the coarse policy, where
\[
 B_{J,M}=2(M+1)\{J(2M+1)+J(J-1)/2\}+2(M+J+1).
\]
For $2\le M$ and $J\le M$, $B_{J,M}\le12JM^2$. Indeed,
$M+1\le3M/2$, $2M+1\le5M/2$, and $J-1\le M$ bound the main term by
$9JM^2$, and the final term by $3JM^2$. The definition of $H$ gives
$B_{J,M}q_f\le V/8$, while $q_c\le q_f$ gives coarse area at most $9V/8$.
Both statements hold on the same fixed-checkpoint coverage event of
probability at least $1-\delta/2$.

For rapid maintenance the area-to-regret lower bound applies. For its upper
bound, every available arm in this family has gap at most $2\varepsilon$.
Hence, on every path, $R_T\le2\varepsilon T$. Since
$(J+1)(M+J+1)/(JM)\le5$,
\[
 2\varepsilon T\le1280\varepsilon JM q_f.
\]
Furthermore $K\le2M^2$, $K+1\le3M^2$ imply
$L_K\le5\log(M/\delta)$, and
$q_f\le165d^{-2}\log(M/\delta)$. Thus
\[
 2^{-20}\frac{\varepsilon JM}{d^2}\log\frac M\delta
 \le\mathfrak R_\Theta(V/8)
 \le2^{18}\frac{\varepsilon JM}{d^2}\log\frac M\delta.
\]
This upper bound uses the stated horizon and must not be extended unchanged
to arbitrarily long windows.

For the coarse policy, within-checkpoint candidate regret is at most
$2\varepsilon JM q_c$. Candidates have
$\gamma_i\ge\varepsilon-d\ge3\varepsilon/4$, $D_i^T\le2\varepsilon$;
old references have $\gamma_i\ge v$, $D_i^T\le Jv$. Dropping the favorable
quota subtraction in the preceding reward bound and absorbing integer terms
gives
\[
 \sup_\Theta ER_T\le2\varepsilon JM q_c+
 64\left(\frac{JM}{\varepsilon}+\frac{J^2}{v}\right)\log(2T)+1.
\]
At fixed $J,M,\varepsilon,v,\delta$, $d\downarrow0$ makes
$T=\Theta(d^{-2})$. Consequently the rapid endpoint is $\Theta(d^{-2})$
and the deferred endpoint is at most $O(\log(1/d))$. No matching lower bound
for the latter endpoint is asserted.

\section{Eventual exact classification and unconditional regret}
\label{app:consistency}

Let epochs be $E_e=\{16^e,\ldots,16^{e+1}-1\}$ and use cumulative targets
$q_e=\lceil4^e(e+1)^2\rceil$. At even rounds serve, by FIFO, any arrived
arm below the current target, including arms formerly reported excluded.
Other rounds use full-history UCB. In epoch $e$ report from the raw intervals
with radius
\[
 b_{i,e}(n)=\sqrt{\frac{\log\{16i^2(e+1)^2n^2/\delta\}}{2n}},
\]
clipped to $[0,1]$ and equal to $[0,1]$ at zero count. Intervals are
recomputed rather than permanently intersected across epochs.

Hoeffding and a union over $i,e,n$ bound the probability of any coverage
failure by $\delta\zeta(2)^3/8<\delta$. This proves all-time reporting
reliability. Suppose arrivals eventually stop with $K<\infty$. For all large
$e$, the current epoch's remaining checkpoint work is at most $Kq_e$, and
its $2Kq_e$ physical service time is smaller than $15\cdot16^e$. Therefore
each arm has at least $q_{e-1}$ observations at the beginning of every
sufficiently late epoch. Uniformly for $t\in E_e$ and actual
$n\in[q_{e-1},t]$, the radius is
$O(\sqrt{\{e+\log i+\log(1/\delta)\}/q_{e-1}})=o(1)$.

For every positive rational $z$, Hoeffding makes
$\sum_nP(|\hat r_i(n)-r_i|>z)<\infty$. Tail union bounds imply only
finitely many deviations of size $z$ almost surely. Taking a countable
intersection over $z$ and the finitely many final arms gives simultaneous
empirical-mean convergence. If
$\min_i|\max_jr_j-r_i-\varepsilon|>0$, then all interval comparisons
are eventually correct and complete, almost surely. This conclusion is not
conditioned on early coverage, and early reporting errors can be corrected.

Through $T\in E_e$, service pulls number at most $Kq_e$, because targets are
cumulative and lifetime counts never decrease. This is
$O(K\sqrt T(1+\log T)^2)$. Apply the unconditional UCB counting bound to
the remaining pulls to obtain
\[
 ER_T=O\{K\sqrt T(1+\log T)^2+\sqrt{KT\log(2T)}\}=o(T)
\]
for each eventually finite instance. Exact boundary cases are not included
in the complete-classification claim.

\section{Worst-schedule area and a sparse but bursty schedule}
\label{app:worst_area}

Suppose $0<\varepsilon\le1/8$, $0<\eta\le\varepsilon/32$,
$0<\delta\le1/16$, and $K\ge2$. Let $\mathfrak A_{T,K}$ be the smallest
area threshold attainable with joint label-and-area probability at least
$1-\delta$, uniformly over reward environments and schedules with $K_T\le K$.
The policy must be reliable on the full model.

\paragraph{Upper bound.}
Use fixed precision $\rho=\eta/4$. Then
\[
 q_i=\left\lceil32\eta^{-2}\log\frac{4i(i+1)}\delta\right\rceil
 \le99\eta^{-2}\log(K/\delta),\qquad i\le K.
\]
The arbitrary-arrival queue bound proves, on coverage,
\[
 \mathcal U_T\le\min\{KT,2K\sum_{i\le K}q_i\}
 \le256\min\{KT,K^2\eta^{-2}\log(K/\delta)\}.
\]
The UCB reward bound and at most $\sum_iq_i$ service pulls also give
\[
 ER_T\le\min\{T,K+4\sqrt{3KT\log(2T)}+1+
                  99K\eta^{-2}\log(K/\delta)\}.
\]

\paragraph{Lower bound.}
At round one release a reference of mean $1/2+\varepsilon$ and
$m=K-1$ candidates of means $1/2\pm2\eta$. Their signs are independent and
specify their relative-goodness labels. The reference is always optimal;
its known distribution gives no sign information. Boundary distances equal
$2\eta>\eta$, and two-way KL is at most $128\eta^2$. Set
\[
 q_* =\frac{m\log(m/\delta)}{131072\eta^2},\qquad
 h=\min\{T,\lceil q_*\rceil-1\}.
\]
Since $\eta\le1/256$, $q_*\ge m\log(m/\delta)/2>1$ and
$h\ge\frac12\min\{T,q_*\}$. Single-window truncation gives a fixed
environment with
\[
 P\{\mathcal U_T\ge mh/8\}\ge1/7>\delta.
\]
Therefore no valid uniform area quantile can be smaller than
$m\min\{T,q_*\}/16$. The comparisons
$m\ge K/2$ and $\log(m/\delta)\ge\frac12\log(K/\delta)$ yield
\[
 \mathfrak A_{T,K}\ge
 2^{-24}\min\{KT,K^2\eta^{-2}\log(K/\delta)\}.
\]
The lower bound uses a static first-round burst; it does not attribute its
statistical logarithm uniquely to dynamic arrivals.

At fixed task parameters, uniformly over schedules satisfying a terminal
count limit $K_T\le k(T)$, sublinear expected regret and sublinear area
envelopes are simultaneously attainable if and only if
$k(T)^2\log k(T)=o(T)$. Necessity is the last lower bound; sufficiency follows
from the displayed upper bounds. This terminal-count-class conclusion
allows the worst schedule to depend on $T$.

\paragraph{A fixed sparse schedule that still defeats a uniform area envelope.}
Set
\[
 s_n=2^{2^n},\qquad
 m_n=\left\lfloor\frac{\sqrt{s_n}}{(\log s_n)^{1/4}}\right\rfloor.
\]
Release one reference at time one and $m_n$ candidates at each $s_n$.
Earlier batch sizes sum to $o(m_n)$, so $K_{s_n}/\sqrt{s_n}\to0$; between
batches the ratio decreases. Thus $K_t=o(\sqrt t)$ on this one schedule.
For each selected batch $n$, give the reference mean $1/2+\varepsilon$,
that batch means $1/2\pm2\eta$, and other candidates mean $1/2$. Let
\[
 h_n=\left\lceil\frac{m_n\log(m_n/\delta)}{131072\eta^2}\right\rceil-1,
 \qquad T_n=s_n+h_n-1.
\]
For large $n$, $h_n=o(s_n)$ and the next batch has not arrived. The old
history is independent of the new signs. Single-window truncation supplies
one fixed configuration for this endpoint with
$P(\mathcal U_{T_n}\ge m_nh_n/8)\ge1/7$. But
\[
 \frac{m_nh_n}{T_n}\asymp_{\eta,\delta}\sqrt{\log s_n}\longrightarrow\infty.
\]
No deterministic envelope $a_T=o(T)$ can consequently hold uniformly over
reward environments for this schedule. The reward environment may vary
with the chosen endpoint; this is exactly the uniform-in-environment
statement and is not a claim of simultaneous failure at all endpoints in
one fixed reward environment.
